\ifdefined\pdfsuppressptexinfo
  \pdfsuppressptexinfo=15
\fi
\documentclass[10pt]{article}
\pdfinfoomitdate=1
\pdftrailerid{}

\usepackage[T1]{fontenc}
\usepackage[utf8]{inputenc}
\usepackage{lmodern}
\usepackage[margin=0.9in]{geometry}
\usepackage{amsmath,amssymb,amsthm,mathtools}
\usepackage{booktabs}
\usepackage{array}
\usepackage{enumitem}
\usepackage{microtype}
\usepackage{xcolor}
\usepackage{hyperref}
\makeatletter\g@addto@macro{\UrlBreaks}{\do\_}\makeatother

\hypersetup{
  colorlinks=true,
  linkcolor=blue!55!black,
  citecolor=blue!55!black,
  urlcolor=blue!55!black,
  pdftitle={Phase-erased arrows of time and entropy production under periodic driving},
  pdfauthor={Luqman Mahmood}
}

\newtheorem{theorem}{Theorem}[section]
\newtheorem{proposition}[theorem]{Proposition}
\newtheorem{lemma}[theorem]{Lemma}
\newtheorem{corollary}[theorem]{Corollary}
\theoremstyle{definition}
\newtheorem{definition}[theorem]{Definition}
\newtheorem{remark}[theorem]{Remark}
\newtheorem{example}[theorem]{Example}

\DeclareMathOperator{\KL}{KL}
\DeclareMathOperator{\TV}{TV}
\DeclareMathOperator{\Bern}{Bern}
\newcommand{\per}{\mathcal T}
\newcommand{\sbar}{\bar\sigma}
\newcommand{\Bbar}{\bar B}
\newcommand{\Pbar}{\bar P}
\newcommand{\Rrec}{R'}
\newcommand{\E}{\mathbb E}
\newcommand{\Rr}{\mathbb R}
\newcommand{\Z}{\mathbb Z}
\newcolumntype{L}[1]{>{\raggedright\arraybackslash}p{#1}}

\title{Phase-erased arrows of time and entropy production\\
under periodic driving:\\
a bound under known time symmetry, a no-go without it,\\
and the locality of the pairwise estimator}
\author{Luqman Mahmood}
\date{1 October 2026\\[3pt]\small Preprint\\
DOI: \href{https://doi.org/10.5281/zenodo.23049617}{10.5281/zenodo.23049617}}

\begin{document}

\maketitle

\begin{abstract}
We ask when forward records with the drive phase erased can bound entropy
production in periodically driven finite-state Markov jump processes.
Known time symmetry and a reversal-covariant observation give a
finite-window relative-entropy bound with a mirror-state boundary
allowance; the corresponding rate bound has no boundary term.
The allowance vanishes exactly when all frozen generators share the
periodic state's constant stationary law, a weaker condition than detailed
balance. Without time symmetry, a two-state construction fixes the
four-sample arrow at any prescribed value in an explicit interval while
the window's expected entropy production tends to zero. This excludes a
universal quantitative lower bound from that arrow alone, not qualitative
or model-specific inference. Exact probability tables and numerical
checks accompany the construction. The two-state witness's pairwise
waiting-time estimator is identically zero at every period. A separate
locality result covers graphs with hidden edges under explicit uniformity
hypotheses; it does not prove Maier and collaborators' finite-period
pairwise-estimator conjecture.
The ingredients of the time-symmetric bound are classical; no priority
claim is made.
\end{abstract}

\section{Introduction}

When does an arrow visible after forgetting drive phase bound total entropy
production? The distinction between finite-window boundary terms and long-time
rates is essential, as is the distinction between reversal of a recorded
trajectory and execution of the reversed drive protocol.
The following conditions let a forward-only, phase-erased arrow
bound total entropy production; the
conditions are sufficient (Theorem~T1), not necessary for every
forward-only inference:
\begin{enumerate}[label=(\arabic*),nosep]
\item the protocol is known to be time-symmetric up to a phase shift,
  $Q(\theta)=Q(c-\theta)$;
\item the observation map is reversal-covariant, $O\circ R=\Rrec\circ O$;
\item finite windows are charged the boundary allowance $\Bbar$.
\end{enumerate}
Without known time symmetry no universal certificate
$L\sbar\geq F(\KL^{\mathrm{obs}}_L)$ from this arrow alone exists
(Theorem~T2).  A universal \emph{quantitative} lower bound from the
arrow alone then needs information beyond the arrow; the routes available
here are known time symmetry, reversed-protocol data, or phase-resolved
occupations together with transition rates. Qualitative and model-specific
forward-only inferences are not excluded: in the two-state class of T2 a
positive arrow still implies $\sbar>0$
(Remark~\ref{rem:qualitative}).

\paragraph{Classical baseline.}
The ingredients of Theorem~T1 are a forward/reversed-protocol
likelihood identity and information contraction. Related classical
results are Crooks's fluctuation theorem \cite{Crooks1999},
the Hamiltonian phase-space dissipation identity of Kawai, Parrondo and
Van den Broeck \cite{KPV2007}, and the coarse-graining bounds of
Gomez-Marin, Parrondo and Van den Broeck \cite{GPV2008}.
The finite-state path identity needed here is derived explicitly below.

\paragraph{Verification material.}
The supplement \cite{EP1Brecord} supplies exact probability tables,
numerical checks, source code and a history note identifying invalid
earlier runs. Section~\ref{sec:verification} separates these checks from
the analytic arguments. An earlier version used the three-sample
comparator in a weaker limiting argument; the theorem below uses the
reversal-covariant four-sample arrow and an exact intermediate-value
construction.

\section{Setting}
\label{sec:setting}

\subsection{Periodically driven jump processes, phase erasure, observation maps}

\begin{definition}[Driven process, window entropy production]
\label{def:process}
Let $S$ be a finite set of states, each even under time reversal, and let
$Q(\theta)$, $\theta\in\Rr/\per\Z$, be a $\per$-periodic family of
generators ($Q_{ij}(\theta)\geq0$ for $i\neq j$, rows summing to zero),
with bounded, right-continuous, piecewise continuously differentiable rates
and finitely many switching phases per period,
irreducible, with $Q_{ij}(\theta)>0$ if and only if $Q_{ji}(\theta)>0$, and
with a unique periodic state $p(\theta)$ solving the forward equation
$\dot p=pQ$ almost everywhere with $p(\theta)>0$ componentwise.
Derivatives at switches are interpreted one-sidedly; finite entropy
comparisons assume integrability of the displayed rate-ratio terms. For a window
$[\theta_0,\theta_0+L]$ let $P_{\theta_0}$ be the law of the path
$(x_t)_{0\leq t\leq L}$ of the process started in $p(\theta_0)$ and driven
by $Q(\theta_0+t)$.  Crooks's total entropy production of the window is
\begin{equation}
\label{eq:sigma-tot}
\Sigma_{\mathrm{tot}}(L;\theta_0)
=\ln p_{\theta_0}(x_0)-\ln p_{\theta_0+L}(x_L)
+\sum_{\text{jumps }x_-\to x_+\text{ at }t}
\ln\frac{Q_{x_-x_+}(\theta_0+t)}{Q_{x_+x_-}(\theta_0+t)},
\end{equation}
the system term plus the medium term.  Writing $\sigma(\theta)=
\sum_{i<j}\bigl(p_iQ_{ij}-p_jQ_{ji}\bigr)\ln\frac{p_iQ_{ij}}{p_jQ_{ji}}(\theta)
\geq0$ for the total entropy production rate in the periodic state, one
has $\E\,\Sigma_{\mathrm{tot}}(L;\theta_0)=\int_0^L\sigma(\theta_0+t)\,dt$,
and the period-averaged rate is $\sbar=\per^{-1}\int_0^\per\sigma(\theta)\,d\theta$.
Since the system entropy is periodic, $\sbar$ is also the period-averaged
medium entropy production rate.
\end{definition}

\begin{definition}[Phase erasure, covariant observation, observed arrow]
\label{def:erasure}
The starting phase $\theta_0$ is uniform on $[0,\per)$ and independent of
the noise; the phase-erased path law is
$\Pbar_L=\per^{-1}\int_0^\per P_{\theta_0}\,d\theta_0$.  Let $R$ be path
reversal on the window, $(Rx)_t=x_{L-t}$.  An observation map $O$ from
paths to records is \emph{phase-free} if it does not use $\theta_0$ and
\emph{reversal-covariant} if there is an involution $\Rrec$ of records with
$O\circ R=\Rrec\circ O$.  Examples: the state path itself ($\Rrec$ = path
reversal); the states sampled on a grid symmetric under $t\mapsto L-t$
($\Rrec$ = order reversal); memoryless state-dependent emissions with
time-even marks; transition records with a declared mark involution.
Detector memory, first-event selection and dead time can break covariance;
they require a separate check of the observation map. The observed arrow and the
full arrow are
\begin{equation}
\label{eq:arrows}
\KL^{\mathrm{obs}}_L:=D\bigl(\Pbar^{\mathrm{obs}}_L\,\big\|\,\Rrec_*\Pbar^{\mathrm{obs}}_L\bigr),
\qquad
\KL^{\mathrm{full}}_L:=D\bigl(\Pbar_L\,\big\|\,R_*\Pbar_L\bigr),
\qquad \Pbar^{\mathrm{obs}}_L=O_*\Pbar_L ,
\end{equation}
where $D(P\|Q)=\int\ln\frac{dP}{dQ}\,dP$ is relative entropy.  Both are
computed from the forward experiment alone: no reversed protocol is run.
\end{definition}

\begin{definition}[Time symmetry, boundary allowance]
\label{def:symmetry}
The protocol is \emph{time-symmetric} if there is a $c$ with
$Q(\theta)=Q(c-\theta)$ almost everywhere. Equalities between protocols
are understood in this sense; values at individual switching times do not
change the path law. For such $c$ the boundary
allowance is
\begin{equation}
\label{eq:Bbar}
\Bbar:=\frac1\per\int_0^\per D\bigl(p(\psi)\,\big\|\,p(c-\psi)\bigr)\,d\psi\;\geq0,
\end{equation}
which is finite when $p>0$.  It measures how far the periodic state is
from its mirror image.
\end{definition}

\subsection{The two-state model}
\label{sec:two-state}

\begin{definition}[Two-state family]
\label{def:two-state}
Take $S=\{0,1\}$, $Q_{01}(t)=u(t)$, $Q_{10}(t)=1-u(t)$, with $u$
$\per$-periodic and $0<u<1$; the total rate is $1$.  Let $p(t)$ be the periodic probability of state~$1$;
it solves $\dot p=u-p$, so $p(t)=\int_0^\infty e^{-r}u(t-r)\,dr$.  For the
trigonometric drive
\begin{equation}
\label{eq:drive}
u(t)=m+e_1\cos\omega t+e_2\cos(2\omega t+\varphi),\qquad \omega=2\pi/\per,
\end{equation}
the periodic state is in closed form,
\begin{equation}
\label{eq:pstate}
p(t)=m+e_1A_1\cos(\omega t-\alpha_1)+e_2A_2\cos(2\omega t+\varphi-\alpha_2),
\qquad A_k=\bigl(1+(k\omega)^2\bigr)^{-1/2},\quad \alpha_k=\arctan(k\omega).
\end{equation}
The entropy production rate is
\begin{equation}
\label{eq:ep-rate}
\sigma(t)=\bigl(u(t)-p(t)\bigr)\ln\frac{u(t)\,(1-p(t))}{(1-u(t))\,p(t)}\;\geq0,
\end{equation}
and $\per\sbar=\int_0^\per\sigma(t)\,dt$ is the entropy production per
period.  For $e_1e_2\neq0$ the drive \eqref{eq:drive} is time-symmetric if
and only if $\varphi\equiv0\pmod\pi$, and then $c=0$.  With $c=0$ the
allowance \eqref{eq:Bbar} reads
$\Bbar=\per^{-1}\int_0^\per\bigl[p\ln\frac{p}{p^-}+(1-p)\ln\frac{1-p}{1-p^-}\bigr](t)\,dt$
with $p^-(t)=p(-t)$.
\end{definition}

\begin{definition}[Phase-erased three-sample comparator]
\label{def:kl3}
The propagator of the two-state process is
$\Pr(X_{t_1}=1\mid X_{t_0}=x_0)=p(t_1)+\bigl(x_0-p(t_0)\bigr)e^{-(t_1-t_0)}$.
For lags $a,b>0$ let $\Pbar_{a,b}(x_0,x_1,x_2)$ be the phase-erased law of
$\bigl(X_{\theta_0},X_{\theta_0+a},X_{\theta_0+a+b}\bigr)$, that is
\begin{equation}
\label{eq:triple}
\Pbar_{a,b}(x_0,x_1,x_2)=\frac1\per\int_0^\per \pi_{x_0}(t)\,
G(x_1\mid x_0;t,t+a)\,G(x_2\mid x_1;t+a,t+a+b)\,dt,
\end{equation}
with $\pi_1=p$, $\pi_0=1-p$ and $G$ the propagator.  The
\emph{phase-erased three-sample comparator} is
\begin{equation}
\label{eq:kl3}
\KL_3(a,b):=\sum_{x\in\{0,1\}^3}\Pbar_{a,b}(x_0,x_1,x_2)\,
\ln\frac{\Pbar_{a,b}(x_0,x_1,x_2)}{\Pbar_{b,a}(x_2,x_1,x_0)} .
\end{equation}
Its quasi-static counterpart at phase lags $\alpha,\beta$ is
\begin{equation}
\label{eq:cstar}
c_*(\alpha,\beta):=\sum_{x}P^{\mathrm{qs}}_{\alpha,\beta}(x_0,x_1,x_2)\,
\ln\frac{P^{\mathrm{qs}}_{\alpha,\beta}(x_0,x_1,x_2)}{P^{\mathrm{qs}}_{\beta,\alpha}(x_2,x_1,x_0)},
\qquad
P^{\mathrm{qs}}_{\alpha,\beta}(x)=\frac1{2\pi}\int_0^{2\pi}
\prod_{k=0}^{2}\Bern\bigl(u(\theta+\tau_k),x_k\bigr)\,d\theta,
\end{equation}
where $(\tau_0,\tau_1,\tau_2)=(0,\alpha,\alpha+\beta)$, $u(\theta)=m+e_1\cos\theta+e_2\cos(2\theta+\varphi)$
in phase units, and $\Bern(q,1)=q$, $\Bern(q,0)=1-q$.  We also write
$c_3^*:=c_*$.
\end{definition}

\begin{definition}[Covariant four-sample arrow]
\label{def:kl4}
For lags $a,b>0$ and $L=a+b$ let $O_4$ record the states on the
reflection-symmetric grid $\{0,a,b,L\}$ (in increasing order), with
$\Rrec$ the order reversal of records.  The \emph{covariant four-sample
arrow} is its observed arrow in the sense of
Definition~\ref{def:erasure},
\[
\KL_4(a,b):=\KL^{\mathrm{obs}}_L
=D\bigl(O_{4*}\Pbar_L\,\big\|\,\Rrec_*O_{4*}\Pbar_L\bigr),
\]
and its quasi-static counterpart at phase lags $\alpha,\beta$,
$c_4^*(\alpha,\beta)$, is defined as in \eqref{eq:cstar} with the four
phases $(0,\alpha,\beta,\alpha+\beta)$ and the reversed record in the
denominator.
\end{definition}

\begin{lemma}[$\KL_3$ is dominated by the covariant arrow]
\label{lem:kl3-observed}
Let $L=a+b$ and let $O_3$ sample the state at times $0$, $a$ and $a+b$.
Then
\[
\KL_3(a,b)=D\bigl(O_{3*}\Pbar_L\,\|\,O_{3*}R_*\Pbar_L\bigr)\;\leq\;\KL_4(a,b)\;\leq\;\KL^{\mathrm{full}}_L .
\]
For $a\neq b$ the map $O_3$ is not reversal-covariant, so $\KL_3$ is a
forward-only comparator (both of its laws are forward phase-erased laws,
at lag orders $(a,b)$ and $(b,a)$), not an observed arrow in the sense of
Definition~\ref{def:erasure}; it is the $(0,a,L)$-marginal comparison of
the covariant arrow $\KL_4$.
\end{lemma}

\begin{proof}
The reversed path sampled at $0$, $a$, $a+b$ is
$(O_3\circ R)(x)=(x_L,x_{L-a},x_0)=(x_{a+b},x_b,x_0)$, which uses the
unrecorded state $x_b$.  Its law under $\Pbar_L$, evaluated at
$(x_0,x_1,x_2)$, is $\Pbar_{b,a}(x_2,x_1,x_0)$, the denominator in
\eqref{eq:kl3}.  For the grid $\{0,a,b,a+b\}$, reversal maps
$(x_0,x_a,x_b,x_L)$ to $(x_L,x_b,x_a,x_0)$, an order reversal since
$L-a=b$; the marginal of $\Rrec_*O_{4*}\Pbar_L$ on the coordinates
$(0,a,L)$ is again $\Pbar_{b,a}(x_2,x_1,x_0)$, and the marginal of
$O_{4*}\Pbar_L$ is $\Pbar_{a,b}$.  Data processing for this marginal gives
the first inequality, and for $O_4$ the second.
\end{proof}

\begin{remark}[Strict inequality at the witness]
\label{rem:strict}
At the witness of Section~\ref{sec:witness} the quasi-static
four-sample law, $(m,e_1,e_2,\varphi)=(9/20,1/4,1/5,\pi/3)$ on the phases
$(0,\pi/3,2\pi/3,\pi)$, has sixteen positive rational cells with
denominator $640000$ (Appendix~\ref{app:add1}).  The records $0000$ and
$0010$ both project to the triple $000$.  Record $0000$ has forward and
reversed probabilities $50681/640000$ and $50681/640000$; record $0010$
has $42799/640000$ and $45799/640000$.  So the likelihood ratio takes the
two values $1$ and $42799/45799$ on one fibre.  By the equality case of data processing,
$c_4^*>c_3^*$ strictly; in fact every one of the eight fibres is
non-constant. Two computational routes reproduce this exact witness.  The values are
\[
c_3^*=0.000738887716627336479919968717667\ldots,\qquad
c_4^*=0.00153441999418270944116752165504\ldots,
\]
agreeing to $45$ digits on the two routes. The independently evaluated
30-digit quadrature result differs by relative error $1.5\cdot10^{-28}$.
\end{remark}

\begin{remark}[Moment structure of the quasi-static comparator]
\label{rem:moments}
The eight probabilities $P^{\mathrm{qs}}_{\alpha,\beta}(x)$ are determined by
the moments $\E u_k$, $\E u_ku_l$, $\E u_0u_1u_2$ of $u_k=u(\theta+\tau_k)$
under uniform $\theta$.  The single moments are all $m$; the pair moments
are values of the autocorrelation
$\E u(\theta)u(\theta+\tau)=m^2+\tfrac12e_1^2\cos\tau+\tfrac12e_2^2\cos2\tau$,
even in $\tau$, at the same three lags $\alpha,\beta,\alpha+\beta$ for the
forward and the reversed law.  The two laws can differ only through the
triple moment, $\E\,u(\theta)u(\theta+\alpha)u(\theta+\alpha+\beta)$
against $\E\,u(\theta)u(\theta+\beta)u(\theta+\alpha+\beta)$, and within
it only through the (1,1,2)-harmonic term proportional to $e_1^2e_2$, the
only one carrying the biphase $\varphi$.  At $\varphi\equiv0\pmod\pi$ the
two agree (time symmetry); at the lags $(\alpha,\beta)=(\pi/2,\pi)$ they
agree identically in $\varphi$. This (1,2)-harmonic cancellation gives
$c_*=0$ at those lags even for $\varphi=\pi/3$, providing a zero control
for the numerical checks.
\end{remark}

\section{Theorem T1: known time symmetry survives phase erasure}
\label{sec:T1}

\begin{theorem}[T1]
\label{thm:T1}
Let $Q$ be time-symmetric, $Q(\theta)=Q(c-\theta)$, with periodic state
$p>0$, and let $O$ be phase-free and reversal-covariant.  Then for every
window length $L>0$,
\begin{equation}
\label{eq:T1}
\KL^{\mathrm{obs}}_L\;\leq\;L\,\sbar+\Bbar .
\end{equation}
\end{theorem}

\begin{proof}
\emph{Step 1 (the reversed experiment is a forward experiment).}  The
reversed protocol of the window $[\theta_0,\theta_0+L]$ is
$\tilde Q(t)=Q(\theta_0+L-t)$, $0\leq t\leq L$.  By symmetry,
$Q(\theta_0+L-t)=Q(c-\theta_0-L+t)$, so the process under the reversed
protocol, started in the periodic state at its initial phase, is the
forward process started at phase $c-\theta_0-L$; its law is
$P_{c-\theta_0-L}$.

\emph{Step 2 (Radon--Nikodym derivative).}  Against the reversed paths of
that process, the log-likelihood ratio of a forward path is
$\ln p_{\theta_0}(x_0)-\ln p_{c-\theta_0-L}(x_L)$ plus the medium term of
\eqref{eq:sigma-tot}, the rates $Q_{ij}(\theta)>0\Leftrightarrow Q_{ji}(\theta)>0$
making it finite.  Adding and subtracting $\ln p_{\theta_0+L}(x_L)$ and
taking the $P_{\theta_0}$-expectation,
\begin{equation}
\label{eq:step2}
D\bigl(P_{\theta_0}\,\big\|\,R_*P_{c-\theta_0-L}\bigr)
=\E\,\Sigma_{\mathrm{tot}}(L;\theta_0)
+D\bigl(p(\theta_0+L)\,\big\|\,p(c-\theta_0-L)\bigr),
\end{equation}
because $x_L$ has law $p(\theta_0+L)$ under $P_{\theta_0}$.

\emph{Step 3 (phase average).}  Average \eqref{eq:step2} over uniform
$\theta_0$.  On the right, $\E_{\theta_0}\E\,\Sigma_{\mathrm{tot}}(L;\theta_0)
=\E_{\theta_0}\int_0^L\sigma(\theta_0+t)\,dt=L\sbar$ (the system term
averages to zero by periodicity), and the substitution
$\psi=\theta_0+L$ turns the second term into $\Bbar$.  On the left, the
map $\theta_0\mapsto c-\theta_0-L$ preserves the uniform law on the
circle, so $\per^{-1}\int R_*P_{c-\theta_0-L}\,d\theta_0=R_*\Pbar_L$, and
joint convexity of relative entropy gives
$\KL^{\mathrm{full}}_L=D(\Pbar_L\|R_*\Pbar_L)\leq L\sbar+\Bbar$.

\emph{Step 4 (data processing).}  $D(O_*\Pbar_L\|O_*R_*\Pbar_L)\leq
D(\Pbar_L\|R_*\Pbar_L)$, and by covariance $O_*R_*\Pbar_L=\Rrec_*O_*\Pbar_L
=\Rrec_*\Pbar^{\mathrm{obs}}_L$.  Since $O$ is phase-free it commutes with
the phase mixture, so the left side is $\KL^{\mathrm{obs}}_L$.
\end{proof}

\begin{corollary}[Phase loss and observation loss under known symmetry]
\label{cor:information-loss}
Under the hypotheses of Theorem~\ref{thm:T1}, assume the relative entropies
below are finite.  On phase--path space define
\[
 \mathbb P(d\theta,dx)=\frac{d\theta}{\per}P_\theta(dx),\qquad
 \mathbb Q(d\theta,dx)=\frac{d\theta}{\per}
                 (R_*P_{c-\theta-L})(dx).
\]
Let $P_X,Q_X$ be their path marginals and $P_Y,Q_Y$ the marginals of
$Y=O(X)$.  Then
\begin{equation}\label{eq:information-loss}
 L\sbar+\Bbar
 =\KL^{\mathrm{obs}}_L+I_{\mathrm{phase}}+I_{\mathrm{obs}},
\end{equation}
where the two nonnegative terms are
\begin{align*}
 I_{\mathrm{phase}}
 &=\int D\bigl(\mathbb P(d\theta\mid x)\,
                  \big\|\,\mathbb Q(d\theta\mid x)\bigr)\,P_X(dx),\\
 I_{\mathrm{obs}}
 &=\int D\bigl(P_X(dx\mid y)\,\big\|\,Q_X(dx\mid y)\bigr)\,P_Y(dy).
\end{align*}
These quantify the losses from first erasing phase and then reducing the
path to the recorded observable, in that order.
\end{corollary}
\begin{proof}
The phase marginals of $\mathbb P$ and $\mathbb Q$ are the same uniform
law.  Steps 2--3 of the proof of Theorem~\ref{thm:T1} therefore give
$D(\mathbb P\|\mathbb Q)=L\sbar+\Bbar$.  The relative-entropy chain rule,
first conditioned on $X$ and then on $Y=O(X)$, gives
\[
 D(\mathbb P\|\mathbb Q)=D(P_X\|Q_X)+I_{\mathrm{phase}}
 =D(P_Y\|Q_Y)+I_{\mathrm{obs}}+I_{\mathrm{phase}}.
\]
The phase shift in $\mathbb Q$ preserves the uniform law, so
$Q_X=R_*\Pbar_L$; reversal covariance identifies
$D(P_Y\|Q_Y)=\KL^{\mathrm{obs}}_L$.
\end{proof}
This is a chain-rule consequence in the known-symmetry setting, not a new
universal entropy identity.  In particular the reversed joint law uses
$c-\theta-L$, not the unshifted phase.  The finite-window allowance
$\Bbar$ remains part of the left side.  The two losses depend on this
order of coarse-graining and do not assert two independent physical
mechanisms.  The decomposition does not extend the bound to the
unknown-symmetry counterexample.

\begin{corollary}[rate form and score form]
\label{cor:forms}
Under the hypotheses of Theorem~\ref{thm:T1}:
\begin{enumerate}[label=(\alph*),nosep]
\item (rate form) $\displaystyle \sbar\;\geq\;\limsup_{L\to\infty}\frac{\KL^{\mathrm{obs}}_L}{L}$,
  with no boundary term;
\item (score form) for any record score $S$ with values in $[-1,1]$ that is
  odd under $\Rrec$, $S\circ\Rrec=-S$, with phase-erased mean
  $\mu_L=\E_{\Pbar^{\mathrm{obs}}_L}S$,
  \begin{equation}
  \label{eq:score}
  L\,\sbar\;\geq\;2\mu_L^2-\Bbar .
  \end{equation}
\end{enumerate}
\end{corollary}

\begin{proof}
(a) Divide \eqref{eq:T1} by $L$; $\Bbar$ is finite and independent of $L$.
(b) $\E_{\Pbar^{\mathrm{obs}}}S-\E_{\Rrec_*\Pbar^{\mathrm{obs}}}S
=\mu_L-\E_{\Pbar^{\mathrm{obs}}}(S\circ\Rrec)=2\mu_L$, while
$|\E_PS-\E_QS|\leq2\TV(P,Q)\sup|S|$; hence $\TV\geq|\mu_L|$ and Pinsker's
inequality $D\geq2\TV^2$ gives $2\mu_L^2\leq\KL^{\mathrm{obs}}_L\leq L\sbar+\Bbar$.
\end{proof}

The rate form needs no model-dependent
allowance.  Finite windows must be charged $\Bbar$, which is
model-dependent and bounded by population ratios when $p\geq p_{\min}>0$;
short windows can be dominated by it (Section~\ref{sec:verification}
shows a window where it is comparable to $L\sbar$).

\begin{proposition}[When the allowance vanishes]
\label{prop:Bbar-zero}
Let $Q(\psi)=Q(c-\psi)$ and $p>0$.  Then $\Bbar=0$ if and only if the
periodic state is constant, $p(\psi)\equiv\pi$, with $\pi Q(\psi)=0$ for
every $\psi$; that is, if and only if $p$ is a common stationary law of
every frozen generator.
\end{proposition}

\begin{proof}
$\Bbar=0$ forces $D(p(\psi)\|p(c-\psi))=0$ for almost every $\psi$, hence
$p(\psi)=p(c-\psi)$ for all $\psi$ by continuity. Differentiating almost everywhere,
$\dot p(\psi)=-\dot p(c-\psi)=-p(c-\psi)Q(c-\psi)=-p(\psi)Q(\psi)=-\dot p(\psi)$,
so $p\equiv\pi$ with $\pi Q(\psi)=0$ almost everywhere, and hence
everywhere by right continuity of $Q$.
Conversely a constant $p$ equals its mirror image and $\Bbar=0$.
\end{proof}

\begin{remark}[stationarity is weaker than detailed balance]
\label{rem:stationarity}
Common \emph{detailed balance},
$\pi_iQ_{ij}(\psi)=\pi_jQ_{ji}(\psi)$, and preparation in $\pi$
guarantee same-experiment path reversal for a palindromic schedule:
the reversed path weights agree by detailed balance and the equal integrated
escape rates. This is strictly stronger than
Proposition~\ref{prop:Bbar-zero}. For example, a three-state cycle with
constant clockwise rate $k_+$ and
counterclockwise rate $k_-$, $k_+\neq k_-$, is a (constant, hence
time-symmetric) protocol whose uniform stationary law is common to every
frozen generator, so $\Bbar=0$, while detailed balance fails and
$\sbar>0$; Theorem~\ref{thm:T1} then reads $\KL^{\mathrm{obs}}_L\leq L\sbar$.
\end{remark}

\section{Theorem T2: no-go without time symmetry}
\label{sec:T2}

\begin{theorem}[T2]
\label{thm:T2}
Let $\delta=(2-\sqrt3)/8$ and let $\mathcal M_\delta$ be the class of
two-state periodically driven models of Definition~\ref{def:two-state}
whose rate $u=Q_{01}$ and periodic population $p$ lie in $[\delta,9/10]$
(so every rate and population lies in $[\delta,1-\delta]$), observed
through the covariant four-sample record of Definition~\ref{def:kl4} with
lags $(\per/6,\per/3)$ and $L=\per/2$.  Let $F\colon(0,\infty)\to\Rr$ be
any function, with no continuity or monotonicity assumed.  If $F(k_0)>0$
for some $k_0\in(0,c_4^*)$, where $c_4^*=c_4^*(\pi/3,2\pi/3)$ at the
witness drive, then some model in $\mathcal M_\delta$ without time
symmetry has $\KL^{\mathrm{obs}}_L=k_0$ exactly and $L\sbar<F(k_0)$.
Hence no certificate $L\,\sbar\geq F(\KL^{\mathrm{obs}}_L)$ holds on
$\mathcal M_\delta$ without knowledge of time symmetry for any $F$ with
$F(k)>0$ for all $k>0$.  A universal lower bound therefore needs
information beyond this arrow; the routes available here are known time
symmetry (Theorem~\ref{thm:T1}), reversed-protocol data
\cite{HFB2024}, or phase-resolved occupations and transition rates
\cite[Eq.~19]{Maier2024}. The
same holds for the three-sample comparator $\KL_3$ with $c_3^*=c_*$ in
place of $c_4^*$.
\end{theorem}

\begin{remark}[Why an exact fixed-arrow construction is needed]
\label{rem:F}
A sequence with positive limiting arrow and vanishing dissipation excludes
lower-bound functions bounded away from zero near that limiting arrow.
It does not exclude an arbitrary positive function with no such regularity.
The construction below keeps the arrow exactly fixed, eliminating that
restriction. For $F(k)=k$, the finite-period example in
Table~\ref{tab:BA} already violates the proposed bound.
\end{remark}

\subsection{Exact fixed-arrow construction}
\label{sec:T2fixed}

\begin{proof}[Proof of Theorem~\ref{thm:T2}]
Consider the drives
\[
u_\varphi(\theta)=\tfrac14(1+\cos\theta)+\tfrac15\bigl(1+\cos(2\theta+\varphi)\bigr),
\qquad 0\leq\varphi\leq\pi/3,
\]
that is $(m,e_1,e_2)=(9/20,1/4,1/5)$: the witness at
$\varphi=\pi/3$, time-symmetric at $\varphi=0$.  Let $K_\per(\varphi)$ be
the covariant four-sample arrow $\KL_4(\per/6,\per/3)$ of this drive.

\begin{enumerate}[label=\textbf{Step \arabic*.},leftmargin=*,nosep]
\item \emph{Rates.}  $u_\varphi\leq\frac12+\frac25=\frac9{10}$.  If the
  circular distance from $\theta$ to $\pi$ is at least $\pi/6$, then
  $\cos\theta\geq-\cos(\pi/6)$ and the first term is at least
  $\delta=(2-\sqrt3)/8$.  Otherwise $\theta=\pi+\varepsilon$ with
  $|\varepsilon|<\pi/6$, so $2\theta+\varphi\equiv2\varepsilon+\varphi\in
  (-\pi/3,2\pi/3)$, $\cos(2\theta+\varphi)>-\frac12$ and the second term
  exceeds $\frac1{10}>\delta$.  Hence $\delta\leq u_\varphi\leq\frac9{10}$
  for all $\theta$ and all $\varphi\in[0,\pi/3]$.
\item \emph{Populations.}  $p(t)=\int_0^\infty e^{-r}u(t-r)\,dr$ is an
  average of values of $u$, so $\delta\leq p\leq\frac9{10}$.
\item \emph{Lag.}  $|u-p|\leq\int_0^\infty e^{-r}|u(t)-u(t-r)|\,dr\leq
  \sup|\dot u|\int_0^\infty re^{-r}\,dr\leq(\frac14+\frac25)\omega
  =\frac{13}{20}\,\omega$.
\item \emph{Dissipation.}  By \eqref{eq:ep-rate},
  $\sigma=(u-p)(\operatorname{logit}u-\operatorname{logit}p)$, and on
  $[\delta,\frac9{10}]$ the derivative of $\operatorname{logit}$ is
  $1/(x(1-x))\leq10/\delta$.  Hence $\sigma\leq(10/\delta)(\frac{13}{20})^2\omega^2$
  and, with $L=\per/2$ and $\omega=2\pi/\per$,
  \begin{equation}
  \label{eq:uniform-ep}
  L\,\sbar\;\leq\;\frac{169\,\pi^2}{20\,\delta\,\per}
  \qquad\text{uniformly in }\varphi\in[0,\pi/3].
  \end{equation}
\item \emph{Continuity.}  For fixed $\per$ the sixteen record
  probabilities are phase averages of products of the periodic population
  and transition probabilities of an irreducible two-state chain with
  rates $\geq\delta$; they are positive and continuous in $\varphi$, so
  $K_\per$ is continuous on $[0,\pi/3]$.
\item \emph{Limits.}  As $\per\to\infty$ at fixed $\varphi$,
  $e^{-\per/6}\to0$ and $p\to u$ uniformly (Step~3), so the record law
  converges to the positive quasi-static law and $K_\per(\varphi)$
  converges to the relative entropy between that law and its reversal.
  At $\varphi=0$ the drive is even and the grid is
  reflection-symmetric, so the quasi-static law is reversal-invariant and
  $K_\per(0)\to0$; at $\varphi=\pi/3$, $K_\per(\pi/3)\to c_4^*>0$.
\item \emph{Intermediate values.}  Fix $k_0\in(0,c_4^*)$.  For all large
  $\per$, $K_\per(0)<k_0<K_\per(\pi/3)$, and the intermediate value
  theorem gives $\varphi_\per\in(0,\pi/3)$ with $K_\per(\varphi_\per)=k_0$
  exactly.  That model is not time-symmetric ($\varphi_\per\notin\pi\Z$,
  $e_1e_2\neq0$), and by \eqref{eq:uniform-ep} its $L\sbar\to0$, so
  $L\sbar<F(k_0)$ once $\per$ is large.  This proves
  Theorem~\ref{thm:T2}; the comparator case runs identically on the
  $(0,\pi/3,\pi)$ marginals, whose quasi-static comparison is $c_*$ at
  $\varphi=\pi/3$ and $0$ at $\varphi=0$.
\end{enumerate}
\end{proof}

\noindent The argument is an existence proof: it gives no explicit
$\varphi_\per$ and no calibration prescription, it covers the declared
two-state class only, and prescribed arrows $k_0\geq c_4^*$ are not
covered by this family.  The constant in \eqref{eq:uniform-ep} is loose
(at $\per=40960$ it is $0.0608$ against the computed $L\sbar=0.000264$),
but it tends to zero, which is all the argument needs.  As an
illustration only (one route, not part of the proof), at $\per=40960$ the
exact finite-period law gives $K_\per(0)=1.06\cdot10^{-25}$ and
$K_\per(\pi/3)=0.0015344197795$, so the bracket
$K_\per(0)<c_*/2<K_\per(\pi/3)$ already holds there.

\begin{remark}[Scope of T1]
\label{rem:T1scope}
The hypotheses of Theorem~\ref{thm:T1} (known time symmetry, a phase-free
reversal-covariant observation) are sufficient for the bound.  Their
failure does not forbid every model-specific forward-only inference; T2
concerns a universal lower bound from this arrow alone.
\end{remark}

\begin{remark}[qualitative implications survive T2]
\label{rem:qualitative}
Theorem~\ref{thm:T2} excludes universal \emph{quantitative}
lower bounds positive somewhere in $(0,c_4^*)$; it does not remove qualitative or model-specific
implications.  In $\mathcal M_\delta$ itself, $\sbar=0$ forces
$\sigma\equiv0$, because the integrand \eqref{eq:ep-rate} is $\geq0$;
then $u\equiv p$, since $\operatorname{logit}$ is strictly increasing;
then $\dot p=u-p=0$, so the drive is constant.  The process is then a
stationary time-homogeneous two-state chain, which is reversible, so
$\KL^{\mathrm{obs}}_L=0$.  Hence in $\mathcal M_\delta$ a positive observed
arrow does imply $\sbar>0$. What fails is a universal function with
$F(k_0)>0$ at a prescribed $k_0\in(0,c_4^*)$ bounding $L\sbar$ from below. Whether a qualitative
implication of this kind holds beyond the two-state class is not
addressed here.
\end{remark}

\subsection{A finite-period witness}
\label{sec:witness}

The witness is the drive
\begin{equation}
\label{eq:witness}
u(t)=0.45+0.25\cos\omega t+0.2\cos(2\omega t+\pi/3),
\end{equation}
that is $(m,e_1,e_2,\varphi)=(0.45,0.25,0.2,\pi/3)$, with lags
$(a,b)=(\per/6,\per/3)$, so $L=\per/2$ and $(\alpha,\beta)=(\pi/3,2\pi/3)$,
and $\per\in\{40,160,640,2560,10240,40960\}$.
A phase grid estimates the rate extrema as approximately \(0.0997\)
and \(0.8742\); the analytic positive lower bound used in the theorem
is \(\delta\), not the grid estimate.
Table~\ref{tab:BA} shows the \(1/\per\) decay of entropy production
per period and convergence of the comparator to its positive
quasi-static value. Full-precision outputs are supplied in
\path{supplement/results/EP1B_RUN2.json}.

\begin{table}[ht]
\centering
\caption{The asymmetric two-state witness. The observation window has
length \(L=\per/2\). Values are rounded from the supplied numerical record.}
\label{tab:BA}
\begin{tabular}{@{}rllll@{}}
\toprule
$\per$ & EP per period & $\per\cdot$EP & $\KL_3$, the comparator & $L\sbar$ over the window\\
\midrule
40 & 0.498 & 19.93 & $6.43\cdot10^{-4}$ & 0.249\\
160 & 0.134 & 21.47 & $7.32\cdot10^{-4}$ & 0.0671\\
640 & 0.0337 & 21.60 & $7.385\cdot10^{-4}$ & 0.0169\\
2560 & 0.00844 & 21.608 & $7.3886\cdot10^{-4}$ & 0.00422\\
10240 & 0.00211 & 21.610 & $7.38886\cdot10^{-4}$ & 0.00106\\
40960 & 0.000528 & 21.6105 & $7.38888\cdot10^{-4}$ & \textbf{0.000264}\\
\bottomrule
\end{tabular}
\end{table}

At \(\per=40960\), the comparator is approximately
\(0.000738887612\), whereas \(L\bar\sigma\) is approximately
\(0.000263800024\). Thus \(\KL_4\geq\KL_3>L\bar\sigma\).
The positive limiting arrow describes the drive's time asymmetry,
even as dissipation per period tends to zero. Every frozen generator
of this two-state model satisfies detailed balance; the period-scale
arrow does not establish broken detailed balance at a frozen phase.

\section{Local waiting-time statistics in the slow-driving limit}
\label{sec:T3}

Maier, Deg\"unther, van der Meer and Seifert \cite{Maier2024} distinguish
three estimators. Their reversed-protocol waiting-time bound is proved
(journal Eq.~16); their forward-only phase-averaged estimator is
conjectured to be a bound (Eqs.~17--18); and their occupation-and-rate
bound is proved without reversing the protocol (Eq.~19).
These are Eqs.~13--16, respectively, in arXiv:2406.12657v1.
The forward-only waiting-time expression is
\[
\mathcal E_{\per}=\sum_{I,J} n_I^\per
\int_0^\infty \Psi^\per_{I\to J}(t)
\log\frac{\Psi^\per_{I\to J}(t)}
{\Psi^\per_{\widetilde J\to\widetilde I}(t)}\,dt,
\]
where \(n_I^\per\) is the phase-averaged frequency of transition \(I\).
Here \(\Psi^\per_{I\to J}\) is the conditional density of the next
observed event and its waiting time, averaged over the starting phase
with event weights \(p_i(\theta)Q_{ij}(\theta)/(\per n_I^\per)\),
not with uniform phase weights.
The next proposition concerns this expression's limit, not the
conjectured finite-period inequality.

\begin{proposition}[Locality under uniform slow driving]
\label{prop:T3}
Let \(Q_\per(t)=q(t/\per)\), where \(q\) is a continuously differentiable,
period-one generator on a finite fixed connected undirected graph.
Assume both directional rates on every edge lie in a fixed interval
\([a,b]\), \(0<a\leq b<\infty\), and every \(q(s)\) is reversible
with stationary law \(\pi(s)\). Observe a fixed nonempty set of directed
edges closed under reversal, with no detector memory or missed events on
those edges. Then, as \(\per\to\infty\),
\[
\mathcal E_{\per}\longrightarrow0,
\qquad
\bar\sigma_\per\longrightarrow0.
\]
For each fixed waiting time \(t>0\), the intensity-weighted pair densities
\(h^\per_{IJ}(t)=n_I^\per\Psi^\per_{I\to J}(t)\) converge to densities
\(h^\infty_{IJ}(t)=h^\infty_{\widetilde J\,\widetilde I}(t)\).
\end{proposition}

\begin{proof}
The fixed graph and uniform rate bounds give a time-uniform minorization
over an interval of fixed length: prescribe a path of at most
\(|S|-1\) edges between any two states, place its jumps in successive
subintervals, and use the lower rate and upper escape-rate bounds.
Consequently the transition operators contract zero-mass signed measures
exponentially, with constants independent of \(\per\).
The stationary law \(\pi(s)\) is continuously differentiable and has
a positive uniform lower bound. Variation of constants for
\(p_\per(t)-\pi(t/\per)\), whose forcing
\(-\pi'(t/\per)/\per\) has zero total mass because \(\sum_i\pi_i(s)=1\), gives
\[
\sup_t\|p_\per(t)-\pi(t/\per)\|=O(\per^{-1}).
\]
Thus the flux differences and their log-ratios in the entropy-production
formula are both \(O(\per^{-1})\), uniformly, and
\(\bar\sigma_\per=O(\per^{-2})\).

Let \(K_s\) have the diagonal of \(q(s)\), its hidden off-diagonal
transitions, and zero off-diagonal entries for observed transitions.
Its semigroup describes evolution killed at the next observed event.
Writing \(U_\per(\theta,\theta+t)\) for the corresponding time-dependent
killed propagator, the exact finite-period density for
\(I=(i,j)\), \(J=(k,l)\) is
\[
h^\per_{IJ}(t)=\frac1\per\int_0^\per
p_i(\theta)Q_{ij}(\theta)
(U_\per(\theta,\theta+t))_{jk}Q_{kl}(\theta+t)\,d\theta.
\]
Continuity of finite-dimensional propagators on every fixed time interval,
uniformly in \(s=\theta/\per\), gives
\[
h^\infty_{IJ}(t)=\int_0^1
 \pi_i(s)q_{ij}(s)(e^{tK_s})_{jk}q_{kl}(s)\,ds .
\]
Detailed balance, including
\(\pi_j(s)(e^{tK_s})_{jk}=\pi_k(s)(e^{tK_s})_{kj}\),
proves the reversal identity for this limit.

It remains to justify integration over arbitrarily long waiting times.
From every state there is a uniformly positive probability of an observed
event during an interval of fixed length, by the same path argument.
The killed propagators therefore have a uniform exponential tail, giving
\(h^\per_{IJ}(t)\leq C e^{-\kappa t}\).
If a pair is reachable through hidden edges, let \(m\) be the shortest
hidden-path length between \(j\) and \(k\). Prescribing just that path
gives \(h^\per_{IJ}(t)\geq c t^m e^{-\Lambda t}\);
the reversed pair has the same support and path length.
These bounds hold for all sufficiently large \(\per\), with constants
depending only on the fixed model. Unsupported pairs are identically
zero in both directions. On supported pairs they imply
\[
\left|h^\per_{IJ}(t)
 \log\frac{h^\per_{IJ}(t)}
 {h^\per_{\widetilde J\,\widetilde I}(t)}\right|
 \leq C' e^{-\kappa t}(1+t+|\log t|),
\]
an integrable bound. Dominated convergence sends the sum of these
intensity-pair relative entropies to zero.

Finally, consecutive-event counts have equal incoming and outgoing
marginals:
\(\sum_J\int_0^\infty h^\per_{IJ}(t)\,dt=n_I^\per\) and
\(\sum_I\int_0^\infty h^\per_{IJ}(t)\,dt=n_J^\per\).
For the incoming identity, every observed event has an almost-surely
unique observed predecessor by the waiting-time tail bound. Decompose
\(p_k(\theta')\) according to the last observed event before \(\theta'\),
multiply by \(Q_{kl}(\theta')\), and average \(\theta'\) over a period.
The estimator is exactly
\[
\mathcal E_\per
=\sum_{I,J}\int_0^\infty h^\per_{IJ}
 \log\frac{h^\per_{IJ}}{h^\per_{\widetilde J\,\widetilde I}}\,dt
-\sum_I n_I^\per\log\frac{n_I^\per}{n_{\widetilde I}^\per}.
\]
Here \(n_I^\per\to\int_0^1\pi_i(s)q_{ij}(s)\,ds\), and these limits
agree for \(I\) and \(\widetilde I\). The second sum also tends to zero.
\end{proof}

The uniform hypotheses make the fixed-waiting-time limit and its
integration explicit. They do not cover rates that vanish, a graph that
changes with the period, or observation windows that grow proportionally
to the period. The last case is precisely the mechanism of
Theorem~\ref{thm:T2}; it does not refute Maier et al.'s distinct
finite-period conjecture.
For the two-state witness itself, the only supported consecutive pairs are
\((0\!\to\!1,1\!\to\!0)\) and \((1\!\to\!0,0\!\to\!1)\).
Both are self-reverse under \((I,J)\mapsto(\widetilde J,\widetilde I)\),
so \(\mathcal E_\per=0\) at every period. The proposition also covers
graphs with hidden edges, where observed pairs need not be self-reverse.

\section{Summary of information requirements}
\label{sec:T4}

\begin{center}
\begin{tabular}{@{}L{4.2cm}L{5.4cm}L{5.4cm}@{}}
\toprule
conclusion & forward-only, phase erased & additional information\\
\midrule
a lower bound on $\sbar$ & yes, from the limiting arrow rate if the protocol is time-symmetric and
$O$ is covariant (T1) & otherwise, for a universal quantitative bound
from this arrow alone (T2); qualitative, model-specific implications
survive (Remark~\ref{rem:qualitative})\\
a finite-window bound & yes, minus the allowance $\Bbar$
(model-dependent; bounded if $p\geq p_{\min}>0$) & --\\
a general periodic protocol & no universal quantitative bound from this
arrow alone (T2) & Harunari--Fiore--Barato's
$\sigma\geq\hat\sigma-X$ uses reversed-protocol data; Maier et al.'s Eq.~19 uses phase-resolved occupations and rates\\
a forward-only pairwise estimator & Maier et al.'s Eq.~18, conjectured;
consistent with Proposition~\ref{prop:T3} & --\\
\bottomrule
\end{tabular}
\end{center}

An observation with memory, event selection or dead time requires its own
reversal-covariance check. The time-symmetric bound does not automatically
apply to an arbitrary detector model.

\section{Computational checks and limitations}
\label{sec:verification}

The supplied numerical route uses Python 3.11 and mpmath 1.3.0 at
30 digits with deterministic quadrature. For the piecewise-constant drive
\(u=1/4,3/4,1/4\) on three successive intervals of length \(h=\log2\),
it reproduces the periodic switching-time populations
\((11/28,9/28,15/28)\), the mean \(3/14\) of \(X(2h)-X(h)\), and
the entropy production \(3\log3/7\) per cycle. The six-period asymmetric
study is in Table~\ref{tab:BA}. A symmetric control with
\((m,e_1,e_2,\varphi)=(0.45,0.25,0.15,\pi)\) checks
\(\KL_3\leq L\bar\sigma+\bar B\) at periods \(10,40,160,640\).
It checks a consequence of the full-path bound, not the four-sample
arrow directly. At period \(10\), its comparator is approximately
\(9.0\cdot10^{-7}\), against \(L\bar\sigma=0.418\) and
\(\bar B=0.083\).

An earlier symmetric control used \(e_2=0.2\), making its rate zero at
phase \(\pi\), outside the declared positive-rate domain. Numerical
evaluation returned infinity and its acceptance test did not reject
a non-finite value. That block is invalid;
its original result is retained and identified in the supplement's
history. The replacement control requires finite outputs and a positive
rate range. The asymmetric values are unchanged.
The numerical infinity is not a divergence theorem: the old drive obeys
\(u(\pi+z)=21z^2/40+O(z^4)\), so for finite period its entropy integrand
has only an integrable logarithmic singularity there. This does not put
the old run inside the positive-rate hypothesis.

\paragraph{Exact probability checks.}
\label{sec:add1check}
The four-sample table in Appendix~\ref{app:add1} is reconstructed by
a Wolfram Engine 14.3 route and a separately coded rational
Laurent-polynomial route over \(\mathbb Q(\sqrt{-3})\), with logarithms
evaluated at 50 digits on the latter. Positivity, normalization, both
marginal identities and the strict relative-entropy inequality are checked.
The two routes agree on \(c_3^*\) and \(c_4^*\) to 45 digits.
A first corroborating grid check used binary floating-point coefficients
and failed at the exact upper endpoint \(9/10\); its failed record and
the exact-coefficient replacement are both supplied.
These high-precision numbers are not interval enclosures. The
continuous bounds and the fixed-arrow existence statement follow from
the written proofs, not from a sampled grid.

The supplement \cite{EP1Brecord} identifies commands, dependencies,
result files and the interpretation of earlier failures. The locality
proposition is analytic, not an output of these numerical checks.
All conclusions concern the stated Markov and observation models; no
physical measurement, detector calibration or priority claim is made.

\section*{Authorship, assistance and rights}
Luqman Mahmood directed this work and is responsible for its claims and final
text. OpenAI Codex and Anthropic Claude assisted with mathematical development,
computation, checking and editing. Internal and automated checks are not
external peer review. Computational tests do not replace proofs of universal
statements, and no physical identification or experiment is implied.
The paper and mathematical data are licensed CC BY 4.0. Software source and
executable mathematical verification scripts are excluded from that grant;
their rights are reserved unless explicitly stated otherwise.

\begin{thebibliography}{99}

\bibitem{Crooks1999}
G. E. Crooks, Entropy production fluctuation theorem and the
nonequilibrium work relation for free energy differences, \emph{Phys.
Rev. E} \textbf{60} (1999), 2721--2726.
\url{https://doi.org/10.1103/PhysRevE.60.2721}.

\bibitem{KPV2007}
R. Kawai, J. M. R. Parrondo and C. Van den Broeck, Dissipation: The
phase-space perspective, \emph{Phys. Rev. Lett.} \textbf{98} (2007),
080602. \url{https://doi.org/10.1103/PhysRevLett.98.080602}.

\bibitem{GPV2008}
A. Gomez-Marin, J. M. R. Parrondo and C. Van den Broeck, Lower bounds on
dissipation upon coarse graining, \emph{Phys. Rev. E} \textbf{78} (2008),
011107. \url{https://doi.org/10.1103/PhysRevE.78.011107}.

\bibitem{Maier2024}
A. M. Maier, J. Deg\"unther, J. van der Meer and U. Seifert, Inferring
kinetics and entropy production from observable transitions in partially
accessible, periodically driven Markov networks, \emph{J. Stat. Phys.}
\textbf{191} (2024), 104.
\url{https://doi.org/10.1007/s10955-024-03315-7};
preprint \url{https://arxiv.org/abs/2406.12657}.

\bibitem{HFB2024}
P. E. Harunari, C. E. Fiore and A. C. Barato, Inference of entropy
production for periodically driven systems, \emph{Phys. Rev. E}
\textbf{110} (2024), 064126.
\url{https://doi.org/10.1103/PhysRevE.110.064126};
preprint \url{https://arxiv.org/abs/2406.12792}.

\bibitem{EP1Brecord}
L. Mahmood, \emph{Phase-erased entropy production: computational supplement},
2026, distributed with this paper; see \path{supplement/README.md} and
\path{supplement/HISTORY.md}. Source and result paths cited above are relative
to that supplement. DOI:
\url{https://doi.org/10.5281/zenodo.23049617}.

\end{thebibliography}

\appendix

\section{Exact quasi-static laws at the witness}
\label{app:add1}

Copied from \texttt{results/EP1B\_ADDENDUM\_1\_CHECK\_R2.json}; the
Wolfram route prints the same rationals.  Drive
$(m,e_1,e_2,\varphi)=(9/20,1/4,1/5,\pi/3)$, phases $(0,\pi/3,2\pi/3,\pi)$.
For each record $x$ the table gives the forward cell $P_4(x)$ and the
reversed cell $R_4(x)=P_4(\text{reverse of }x)$.

\medskip
{\scriptsize\ttfamily
\begin{tabular}{@{}llllll@{}}
\toprule
x & P4(x) & R4(x) & x & P4(x) & R4(x)\\
\midrule
0000 & 50681/640000 & 50681/640000 & 0001 & 51739/640000 & 51739/640000\\
0010 & 42799/640000 & 45799/640000 & 0011 & 51981/640000 & 48981/640000\\
0100 & 45799/640000 & 42799/640000 & 0101 & 28981/640000 & 31981/640000\\
0110 & 47121/640000 & 47121/640000 & 0111 & 32899/640000 & 32899/640000\\
1000 & 51739/640000 & 51739/640000 & 1001 & 43041/640000 & 43041/640000\\
1010 & 31981/640000 & 28981/640000 & 1011 & 28039/640000 & 31039/640000\\
1100 & 48981/640000 & 51981/640000 & 1101 & 31039/640000 & 28039/640000\\
1110 & 32899/640000 & 32899/640000 & 1111 & 20281/640000 & 20281/640000\\
\bottomrule
\end{tabular}

\smallskip
\noindent Triple comparator, coordinates $(0,\pi/3,\pi)$: P3 =
(2337, 2593, 2323, 1547, 2093, 1777, 2047, 1283)/16000 and\\
R3 = (2412, 2518, 2248, 1622, 2018, 1852, 2122, 1208)/16000 on
000, 001, \dots, 111.\\
c3 = 0.000738887716627336479919968717667278764901883127\\
c4 = 0.00153441999418270944116752165504214461953899412\\
c4 - c3 = 0.000795532277555372961247552937375}

\end{document}
